\par\Needspace{6\baselineskip}\medskip
\hypertarget{native-091-article071}{}\pdfbookmark[2]{71.}{article-071}
\begin{center}71.\end{center}
The aberration of the fixed stars, and also that part of the aberration of planets and comets which is due solely to the Earth's motion, arises because the \emph{tube} moves with the Earth as a whole while a ray of light traverses its optical axis. The observed position of a celestial body (also called its apparent position or its position affected by aberration) is determined by the orientation of the optical axis of a telescope placed so that the light ray issuing from the body reaches both ends of that axis along its path. This orientation differs from the true orientation of the light ray in space. Let us distinguish the two instants $t$ and $t'$ at which the ray reaches the front end (the center of the objective lens) and the rear end (the focus of the objective lens), respectively. Let the positions of these ends in space be $a$ and $b$ at the first instant and $a'$ and $b'$ at the second. Clearly the straight line $ab'$ gives the true orientation of the ray in space, whereas the straight line $ab$ or $a'b'$ corresponds to the apparent position (the two lines may be regarded as parallel). It is also readily \SourceLeaf{092}seen that the apparent position does not depend on the length of the tube. The difference in orientation between the straight lines $b'a$ and $ba$ is the aberration that occurs for fixed stars. We pass over the method of calculating it here, taking it as known. For wandering stars, however, this difference is still short of the full aberration. The planet itself changes position while the ray it emitted travels to the Earth, so the orientation of that ray does not correspond to the true geocentric position at the time of observation. Suppose that the light ray entering the tube at time $t$ left the planet at time $T$. Let the planet's position in space at time $T$ be denoted by $P$ and at time $t$ by $p$. Finally, let $A$ be the position of the front end of the tube's axis at time $T$. The following is then clear.
\par $1^{\mathrm{o}}$ The straight line $AP$ gives the true position of the planet at time $T$.
\par $2^{\mathrm{o}}$ The straight line $ap$ gives the true position at time $t$.
\par $3^{\mathrm{o}}$ The straight line $ba$ or $b'a'$ gives the apparent position at time $t$ or $t'$ (the difference between the times may be regarded as infinitely small).
\par $4^{\mathrm{o}}$ The straight line $b'a$ gives the same apparent position freed from the aberration of the fixed stars.

Now the points $P$, $a$, and $b'$ lie on a straight line, and the segments $Pa$ and $ab'$ are proportional to the time intervals $t-T$ and $t'-t$, since light travels at a uniform speed. Owing to the immense speed of light, the interval $t'-T$ is always very small. During it the Earth's motion may be assumed to be along a straight line at a uniform speed. Thus $A$, $a$, and $a'$ also lie on a straight line, and the segments $Aa$ and $aa'$ are likewise proportional to the intervals $t-T$ and $t'-t$. It follows readily that the straight lines $AP$ and $b'a'$ are parallel, and hence that the first and third positions are identical.

The time $t-T$ is the product of the distance $Pa$ and $493^{\mathrm{s}}$, the time light takes to traverse the mean distance of the Earth from the Sun, which we adopted as the unit of distance. In this calculation the distance $Pa$ may also be replaced by $PA$ or $pa$, since the difference can be of no consequence.

From these principles follow three methods of determining the apparent position of a planet or comet for any time $t$. Different circumstances favor different methods.

I. Subtract from the given time the time light takes to travel from the planet to the Earth. This gives the reduced time $T$, for which the true position computed in the usual way will be identical to the apparent position for $t$. To compute the \SourceLeaf{093}time reduction $t-T$, we must know the distance from the Earth. Convenient means of finding it will usually be available, for example from an ephemeris, even one calculated only roughly. Otherwise, a preliminary calculation of the true distance for time $t$ in the usual way, dispensing with excessive precision, will suffice.

II. Compute the true position and the distance for the given time $t$, then derive the time reduction $t-T$ from the distance, and from this, with the aid of the daily motion (in longitude and latitude or in right ascension and declination), reduce the true position to time $T$.

III. Compute the heliocentric position of the Earth for time $t$ and the heliocentric position of the planet for time $T$. Then combine the two in the usual way to obtain the geocentric position of the planet. Adding the aberration of the fixed stars (found by the known method or taken from tables) gives the required apparent position.

The second method, which is commonly used, has the advantage over the others that determining the distance never requires a double calculation. Its disadvantage is that it can be used only if several neighboring positions are calculated or are already known from observations. Otherwise the daily motion could not be taken as given.

The disadvantage of the first and third methods disappears entirely whenever several neighboring positions are to be calculated. Once the distances for some of them are known, the immediately succeeding distances can be inferred very conveniently and with sufficient precision by the usual means. If the distance is known, the first method will generally be preferable to the third because it has no need of the aberration of the fixed stars. If we must resort to a double calculation, the third method has the advantage that at least the Earth's position is retained in the second calculation.

The requirements for the inverse problem now suggest themselves, for example when the true position is to be derived from the apparent position. In method I you leave the position itself unchanged, but convert the time $t$, to which the given position corresponds as an apparent position, to the reduced time $T$, to which the same position will correspond as a true position. In method II you retain the time $t$ and add to the given position the motion during the interval $t-T$, as if reducing the position to time $t+(t-T)$. In method III, after freeing the given position from the aberration of the fixed stars, you \SourceLeaf{094}regard it as the true position for time $T$. The true position of the Earth corresponding to time $t$ must be retained as if it belonged to the reduced time. The usefulness of the third method will become clearer in Book Two.

For completeness, we note that the position of the Sun is affected by aberration just as the position of a planet is. Since both the distance from the Earth and the daily motion are nearly constant, the aberration itself always takes an almost constant value equal to the mean motion of the Sun in $493^{\mathrm{s}}$, hence $=20''{,}25$. Subtract this quantity from the true longitude to obtain the apparent longitude. The exact value of the aberration is proportional to the product of the distance and the daily motion, or equivalently inversely proportional to the distance. Hence this mean value would have to be decreased at apogee by $0''{,}34$ and increased at perigee by the same amount. Our solar tables already include the constant aberration $-20''{,}25$. To obtain the true longitude one must add $20''{,}25$ to the tabulated longitude.

